\section{From physical preparation to the first release}
\label{sec:entry}
The incoming normal equilibrium loses uniform attraction as $p$ tends
to zero. Negative detuning moves the relevant slip transition to
$p\asymp\sqrt\eps$, where normal dissipation and variation of its
energy metric have comparable sizes. We follow the prescribed rigid
trace by coupling a normal energy with its exact weak coefficient.
A first graph correction has identically zero force defect. The
resulting energy estimate controls the outer range by dissipation and
the critical annulus by a bounded Gronwall factor. Although the weak
coefficient loses accuracy there, conversion back to the physical
state multiplies its error by $p^\beta$. This gives the complete
normal state needed for the first release.

The compact classes are those of Section~\ref{sec:model}. In particular,
$r=-\eta=k\sqrt\eps$, $k\in I=[k_-,k_+]\Subset(0,\infty)$,
and all labels, including the original fast triple, are fixed before
$\eps\downarrow0$. Constants below are common on these classes.
Write $\gamma=-g_*>0$ and $\rho=b/p$.

\subsection{The weak coordinate and a finite normal energy}

\paragraph{The rigid first integral.}
Let $\alpha_1=F_{\rm free}p$, $\alpha_2=F_{\rm free}b$ and
$\alpha_3=-Q_{\rm sat}b$, restricted to $y=v=0$ when written as
functions of $(p,b)$. The rigid saturated field is
\begin{equation}
 \begin{gathered}
 \dot p=\alpha_1,\qquad \dot b=\alpha_2+\alpha_3b/p,\\
 F(p,\rho)=\frac{\alpha_2+(\alpha_3-\alpha_1)\rho}{p},\qquad
 \mathscr D=\alpha_1\partial_p+F\partial_\rho.
 \end{gathered}
 \label{ent:rigid-field}
\end{equation}
Here $\alpha_1(G)=-a$, $\alpha_2(G)=-c$,
$\alpha_3(G)=-a\beta$ and $c=a(\beta-1)\ell$.
The scalar equation for $\rho$ has the regular-singular form
\[
 p\rho_p=H(p,\rho),\qquad
 H=\frac{\alpha_2(p,p\rho)+(\alpha_3-\alpha_1)(p,p\rho)\rho}
              {\alpha_1(p,p\rho)},\qquad
 H(0,\rho)=(\beta-1)(\rho+\ell).
\]
To construct its analytic weak curve, put $s=\bar\rho+\ell$.
Its equation is $(p\partial_p-(\beta-1))s=pJ(p,s)$, with $J$
analytic. In absolutely convergent power series the inverse diagonal
operator divides coefficient $n$ by $n-(\beta-1)$, $n\ge1$.
These divisors stay away from zero uniformly for
$2<\beta_-\le\beta\le\beta_+<5/2$. On a sufficiently small common
disk multiplication by $p$ therefore makes the fixed-point equation
for $s$ a contraction. This gives $\bar b=p\bar\rho=-\ell p+O(p^2)$.

For $d_b=\rho-\bar\rho$, subtraction gives exactly
$d_b'=((\beta-1)/p+A_0(p,d_b))d_b$, where $A_0$ is analytic
and bounded: $\partial_\rho H-(\beta-1)$ is divisible by $p$.
Thus $d_b=p^{\beta-1}\zeta$ satisfies
$\zeta'=A_0(p,p^{\beta-1}\zeta)\zeta$, a regular equation down
to $p=0$ on each bounded $\zeta$ set. Define $\mathcal C(p,b)$
as its limiting value there. Its integral equation and variational
equation give
\begin{gather}
 b(p,C)=\bar b(p)+p^\beta\{C+O(p)\},\qquad
 |\nabla_{p,b}\mathcal C|\le Cp^{-\beta},\notag\\
 |b(p,C)-b(p,\widetilde C)|\le Cp^\beta|C-\widetilde C|,
 \label{ent:weak-coordinate}\\
 \rho=-\ell+O(p)+O(p^{\beta-1}),\qquad
 \mathscr D\rho=O(1),\qquad
 \mathscr D^2\rho=O(p^{\beta-3}).
 \label{ent:weak-derivatives}
\end{gather}
For the first line, the variation in $\zeta$ and its inverse are
bounded by $e^{Cp}$; conversion from
$\zeta=(b-\bar b)/p^\beta$ gives the displayed derivatives.
For the second, differentiate the regular equation along the rigid
field. These are values of rigid derivatives at points of a bounded
coefficient wedge, not derivatives of an actual fast transient.
In particular $|F_p|+|F_\rho|\le C/p$ there. The exact first
integral equals the prescribed $C_1$ on the prescribed rigid trace.

\paragraph{A degenerating normal energy.}
The normal decay rate tends to zero with $p$. We therefore seek a
Lyapunov norm whose square decays at a rate proportional to $p$ for
the frozen normal equation in fast time. Converting this norm to
displacement and memory will cost
$p^{-1/2}$, whereas the velocity estimate remains uniform. These
component weights determine how the normal error enters the weak
coordinate estimate below.

Set $N=(Y,V,Z)=(y/\eps^2,v/\eps,z/\eps^2)$ and define
\begin{equation}
 M_o(p)=\begin{pmatrix}0&1&0\\-p&0&p\\1&0&-2\end{pmatrix},\quad
 e_{\rm eq}=(2,0,1)^T,\quad q_N=(-1,0,1)^T,\quad a_1=(-1,2,-1)^T.
 \label{ent:normal-matrix}
\end{equation}
The polynomial of $M_o$ is $z^3+2z^2+pz+p$.
Its simple root $f=-2+p/4+O(p^2)$ is analytic; write the remaining
factor as $z^2-a_cz+d_c$, where
$a_c=-p/4+O(p^2)$ and $d_c=p/2+O(p^2)$. The polynomial projector
\[
 \begin{gathered}
 P_f=\frac{M_o^2-a_cM_o+d_cI}{f^2-a_cf+d_c},\qquad P_c=I-P_f,\\
 c_2=P_c(0,2,-1/2)^T,\quad c_1=(M_o-a_cI)c_2,\quad
 e_f=P_f(0,0,1)^T.
 \end{gathered}
\]
defines an analytic invertible matrix $T_o=(c_1,c_2,e_f)$.
At zero its columns are $e_{\rm eq}$, $(0,2,-1/2)^T$ and $(0,0,1)^T$;
their determinant is $4$, and
\[
 T_o^{-1}M_oT_o=\operatorname{diag}(A_c,f),\qquad
 A_c=\begin{pmatrix}0&1\\-d_c&-h_c\end{pmatrix},\qquad h_c=-a_c.
\]
Put $S_c=\left(\begin{smallmatrix}d_c&h_c/2\\h_c/2&1\end{smallmatrix}\right)$
and $\mathbb S=T_o^{-T}\operatorname{diag}(S_c,1)T_o^{-1}$.
The exact identity $A_c^TS_c+S_cA_c=-h_cS_c$ and $f<-1$ give
\begin{equation}
 \begin{gathered}
 M_o^T\mathbb S+\mathbb SM_o\le-c_0p\mathbb S,\\
 |w_1|+|w_3|+|q_N^Tw|\le Cp^{-1/2}W,\qquad |w_2|\le CW,
 \end{gathered}
 \label{ent:energy-weights}
\end{equation}
where $W=\|w\|_{\mathbb S}$; also
$\|e_{\rm eq}\|_{\mathbb S}\le C\sqrt p$ and
$\|a_1\|_{\mathbb S}+\|(0,1,0)^T\|_{\mathbb S}\le C$.
Indeed $S_c\asymp\operatorname{diag}(p,1)$; the first and third
columns of $T_o$ have second component $O(p)$, giving the stronger
$w_2$ bound. To differentiate the metric, write
$D_o=\operatorname{diag}(S_c,1)$, $\zeta_o=T_o^{-1}w$ and
$B_o=T_o^{-1}T_o'$. Then
\[
 w^T\mathbb S'w=\zeta_o^TD_o'\zeta_o
 -\zeta_o^T(B_o^TD_o+D_oB_o)\zeta_o,
 \qquad |w^T\dot{\mathbb S}w|\le C|\dot p|p^{-1}W^2.
\]
This estimate uses the actual $\dot p$.

\paragraph{Coupled tracking of the normal and weak coordinates.}
The first graph correction is
\begin{equation}
 G_1=\eps a_1F,\qquad N=e_{\rm eq}\rho+G_1+w.
 \label{ent:first-graph}
\end{equation}
It satisfies $M_oG_1=\eps e_{\rm eq}F$ and $q_N^TG_1=0$ exactly.
To keep the finite-profile feedback, evaluate $A_j=\alpha_j$
at the actual mechanical state
$\Xi(p,p\rho,\eps^2(2\rho+G_{1,1}+w_1),\eps(G_{1,2}+w_2))$,
and set
\[
 \begin{gathered}
 d_\eps=\{A_2+(A_3-A_1)\rho\}/p,\qquad
 \mathcal T_v=\mu B(1-\Phi_\theta(u/\eps))\lambda_N,\\
 \mathcal T_b=\mu(1-\Phi_\theta(u/\eps))\lambda_N
                      (\partial_{\dot q}b)M^{-1}Q_T.
 \end{gathered}
\]
With $E=G_1+w$ the exact equations are
\begin{align*}
 \dot p&=A_1,&
 \dot\rho&=d_\eps-(A_3/p)q_N^TE+\mathcal T_b/p,\\
 \eps\dot E&=M_oE-\eps e_{\rm eq}d_\eps+
       \eps(A_3/p)e_{\rm eq}q_N^TE+\mathcal T_ve_V-\eps e_{\rm eq}\mathcal T_b/p,
 &&e_V=(0,1,0)^T.
\end{align*}
Subtracting the derivative of $G_1$ gives, without truncation,
\begin{equation}
 \begin{split}
 \eps\dot w={}&M_ow+\eps\frac{A_3}{p}(e_{\rm eq}+G_{1,\rho})q_N^Tw\\
 &-\eps(e_{\rm eq}+G_{1,\rho})(d_\eps-F)-\eps G_{1,p}(A_1-\alpha_1)
 -\eps^2a_1\mathscr D^2\rho\\
 &\phantom{{}={}}
 +\mathcal T_ve_V-\frac\eps p(e_{\rm eq}+G_{1,\rho})\mathcal T_b.
 \end{split}
 \label{ent:exact-remainder}
\end{equation}
In particular both profile terms remain in the normal equation.

Here are the estimates needed to use this identity. They hold for
$p\ge h\sqrt\eps$, with any fixed $h>0$, on partial intervals with
bounded weak coefficient, strict normal-force bounds,
$W\le\rho_0\sqrt p$, and $a/2\le-\dot p\le2a$.
From \eqref{ent:weak-derivatives},
$G_1=O(\eps)$,
$\|G_{1,p}\|_{\mathbb S}+\|G_{1,\rho}\|_{\mathbb S}\le C\eps/p$,
and $\|e_{\rm eq}+G_{1,\rho}\|_{\mathbb S}\le C_h\sqrt p$.
The last bound uses $\eps/p^{3/2}\le h^{-3/2}\eps^{1/4}$.
Superscript $G$ denotes $w=0$. Since its physical $y,v$ are
$O(\eps^2)$, $|A_j^G-\alpha_j|\le C\eps^2$ and
$|d_\eps^G-F|\le C\eps^2/p$.
The four saturated inhomogeneous terms in
\eqref{ent:exact-remainder} have energy bounds, respectively,
\begin{equation}
 C\eps^2p^{\beta-3},\qquad C\eps^3p^{-1/2},\qquad
 C\eps^4p^{-2},\qquad C\eps^4p^{-1},
 \label{ent:forcing-domination}
\end{equation}
when ordered as $-\eps^2a_1\mathscr D^2\rho$,
$-\eps e_{\rm eq}(d_\eps^G-F)$,
$-\eps G_{1,\rho}(d_\eps^G-F)$ and
$-\eps G_{1,p}(A_1^G-\alpha_1)$.
Relative to $\eps^2p^{\beta-3}$, the last three bounds have ratios
$\eps p^{5/2-\beta}$, $\eps^2p^{1-\beta}$ and
$\eps^2p^{2-\beta}$. For $h\sqrt\eps\le p\le p_i<1$ and the
prescribed compact range $\beta<5/2$, these are uniformly small.
Thus each term is bounded by $C_h\eps^2p^{\beta-3}$.
At $\beta=5/2$ these three ratios are, respectively,
$O(\eps)$, $O_h(\eps^{5/4})$ and $O_h(\eps^{7/4})$ on the same
interval. These bounds describe the endpoint of this forcing
calculation; the proof below uses the prescribed compact range.
The physical differences from this graph are
$(\eps^2w_1,\eps w_2)$, so
$|A_j-A_j^G|\le C\eps W$ and
$|d_\eps-d_\eps^G|\le C\eps W/p$.
The three homogeneous terms have bounds
$C_h\eps W/p$, $C_h\eps^2p^{-1/2}W$ and
$C\eps^3W/p$. A value bound $1-\Phi_\theta(u/\eps)\le CR_\Phi$
gives for the complete profile forcing
$C_hR_\Phi(1+\eps p^{-1/2}+\eps^2p^{-2})\le C_hR_\Phi$.
The metric derivative and \eqref{ent:energy-weights} consequently give
\begin{equation}
 \eps D^+W\le(-c_0p+C_h\eps/p)W
                +C_h\eps^2p^{\beta-3}+C_hR_\Phi.
 \label{ent:signed-energy}
\end{equation}
We retain the positive metric term when crossing the critical annulus.
Finally the rigid field annihilates $\mathcal C$. At equal $(p,b)$,
its difference from the actual field is
\[
 \delta\dot p=A_1-\alpha_1,\quad
 \delta\dot b=(A_2-\alpha_2)+(A_3-\alpha_3)\rho
                         -A_3q_N^Tw+\mathcal T_b,
 \quad |A_j-\alpha_j|\le C(\eps^2+\eps p^{-1/2}W).
\]
Thus \eqref{ent:weak-coordinate} yields the coupled transport bound
\begin{equation}
 \left|\frac{d\mathcal C_\eps}{dp}\right|
 \le C\{\eps^2p^{-\beta}+p^{-\beta-1/2}W+R_\Phi p^{-\beta}\}.
 \label{ent:weak-transport}
\end{equation}

\subsection{Capture and the actual descending threshold}

\begin{lemma}[The actual incoming threshold]\label{lem:detuned-threshold}
The physical initial states reach a common fixed section $p=p_i>0$
in positive-slip contact. On its continuous contact passage there is a
unique first descending crossing $t_e$ of $u(t_e)=\gamma\eps T$,
where $T=m|\log\eps|$ and $m$ is a sufficiently large fixed constant.
If $p_{\rm stop}$ is the zero of the prescribed detuned rigid slip,
then
\begin{gather}
 p_{\rm stop}=ak\sqrt\eps/\gamma+O(\eps),\qquad
 p_e=p_{\rm stop}+a\eps T+R_e,\qquad R_e=o(\eps),
 \label{ent:threshold-slow}\\
 b_e=-\ell p_{\rm stop}+o(\sqrt\eps),\quad
 |X_e-G|=O(\sqrt\eps),\quad
 \lambda_N(t_e)=\ell+O(\sqrt\eps),\quad V_e=O(\eps),\\
 E_{N,e}:=|Y_e+2\ell|+|V_e|+|Z_e+\ell|=O(\sqrt\eps).
 \label{ent:threshold-normal}
\end{gather}
All assertions are uniform on the fixed compact classes.
\end{lemma}
\begin{proof}
For capture, freeze $p>0$, $b<0$ on a closed incoming rigid arc and
put $\ell_f=-b/p$. The frozen fast equations and proper energy are
\[
 \begin{gathered}
 Y'=V,\quad V'=b+pL,\quad Z'=-Z-L,\quad L=(Z-Y)_+,
 \\
 \mathcal E_f=\frac{V^2}{2p}+\frac{Z^2+L^2}{2}+\ell_fY,
 \qquad \mathcal E_f'=-(Z')^2.
 \end{gathered}
\]
Minimization over $Z$ gives potential $\ell_fY$ for $Y\ge0$ and
$Y^2/4+\ell_fY$ for $Y<0$, so sublevels are uniformly compact
on the fixed positive ranges of $p,\ell_f$. In an invariant
zero-dissipation set, $L=-Z$. If $L>0$, then $Y=2Z$, $V=0$,
and $V'=0$ gives $L=\ell_f$; if $L=0$, then $Z=0$, $Y\ge0$,
but $Y''=b<0$ cannot remain there. The invariance principle therefore
gives attraction to $(-2\ell_f,0,-\ell_f)$, including initially
free states. The Hurwitz matrix \eqref{ent:normal-matrix} gives a
common local exponential estimate on this fixed arc. Continuous
dependence and a finite cover of the parameter and initial-state
compact sets give a common finite fast capture time. Over it the
actual slow drift is $O(\eps)$ and the closed-arc slip margin makes
the profile error $O(e^{-c/\eps})$, so the actual trajectory enters
the strict-contact neighborhood too.

Subtract the moving equilibrium there. Its physical-time derivative
is bounded on the fixed arc; the local Lyapunov inequality gives
$|N-N_*(p,b)|\le Ce^{-c(t-t_c)/\eps}+C\eps$.
The time integral of the slow-field discrepancy is $O(\eps)$.
Gronwall and alignment at a transverse section, leaving a fixed
positive time after capture, give
\begin{equation}
 \mathcal C_\eps(p_i)=C_1+O(\eps),\quad W(p_i)=O(\eps),\quad
 t_\eps(p_i)-t_{{\rm rig},\nu}(p_i)=O(\eps).
 \label{ent:fixed-initialization}
\end{equation}
The graph correction is $O(\eps)$ at this section. Capture before it
may contain departures and recontacts.

We now stop on partial intervals by
$|\mathcal C_\eps-C_1|\le1$, $W\le\rho_0\sqrt p$, strict force
bounds, the regular chart and $a/2\le-\dot p\le2a$, adding the
slip bounds below. Choose $p_i,\rho_0$ small, then $K_L$ large,
and finally $\eps$ small. On $p\ge K_L\sqrt\eps$ the positive
term in \eqref{ent:signed-energy} is absorbed, and the rigid slip
slope is $\gamma/a$. Requiring $u\ge c_{\rm sl}p$ gives a profile
tail smaller than every fixed power of $\eps$. Integration in the
actual decreasing $p$ direction gives
\begin{equation}
 W(p)\le H_\eps(p)+C\eps^2p^{\beta-4}+CR_\Phi/p,\qquad
 H_\eps=CW(p_i)e^{-c(p_i^2-p^2)/\eps}.
 \label{ent:safe-tracking}
\end{equation}
Indeed the forcing integral is
$C\eps\int_p^{p_i}e^{-c(s^2-p^2)/\eps}s^{\beta-3}\,ds$;
$s^2-p^2\ge2p(s-p)$ bounds it by $C\eps^2p^{\beta-4}$.
Substitution into \eqref{ent:weak-transport} gives
\begin{equation}
 |\mathcal C_\eps(p)-C_1|\le
 C\{\eps+\eps^2p^{1-\beta}+\eps^2p^{-7/2}
                           +R_\Phi p^{-\beta-1/2}\}.
 \label{ent:coefficient-bound}
\end{equation}
The principal transport integral is
$\eps^2\int_p^{p_i}s^{-9/2}\,ds=O(\eps^2p^{-7/2})$.
The weighted integral of $H_\eps$ is $O(\eps W(p_i))=O(\eps^2)$:
above $p_i/2$ its width is $O(\eps)$, and below $p_i/2$ the fixed
exponential absorbs every endpoint power. Thus the safe endpoint has
$W=O(\eps^{\beta/2})$ and $\mathcal C_\eps-C_1=O(\eps^{1/4})$.
These bounds improve the coefficient and relative-energy stops.
The rigid slip on the bounded weak wedge is at least
$3\gamma p/(4a)$ after reducing $p_i$; its physical normal
correction is $o(p)$ and $k\sqrt\eps/p\le k_+/K_L$.
The large choice of $K_L$ improves the slip stop. The force, chart
and orientation bounds improve too, so the endpoint is reached.

On $h_I\sqrt\eps\le p\le K_L\sqrt\eps$, where
$h_I=ak_-/(2\sup\gamma)>0$, stop additionally at the actual
threshold $u=\gamma\eps T$. Choose
\begin{equation}
 \kappa_{\rm tail}=\alpha\inf\gamma/2,\qquad
 N_{\rm tail}=\kappa_{\rm tail}m>\beta_+/2+3/4.
 \label{ent:tail-choice}
\end{equation}
The coarse energy radius makes the actual and graph slips differ by
$O(\eps)$, so both are at least $(\inf\gamma)\eps T/2$ before
this stop. Hence $R_\Phi\le C\eps^{N_{\rm tail}}$.
Retain the positive $C_h/p$ in the equation for $W$.
Since $\int C_h\,dt/p\le C_h\log(K_L/h_I)$ and
$\Delta t\le C\sqrt\eps$, its bounded Gronwall factor gives
\[
 W\le C\{\eps^{\beta/2}+\eps^{N_{\rm tail}-1/2}\}
       =O(\eps^{\beta/2}).
\]
The three contributions to weak transport on the annulus are
$O(\eps^{(5-\beta)/2})$,
$O(\eps^{1/4})+O(\eps^{N_{\rm tail}-\beta/2-1/4})$, and
$O(\eps^{N_{\rm tail}+(1-\beta)/2})$.
They give $\mathcal C_\eps-C_1=O(\eps^{1/4})$ throughout;
\eqref{ent:coefficient-bound} also holds there up to a fixed factor.
The component weights now yield
$w_1,w_3,q_N^Tw=O(\eps^{\beta/2-1/4})=o(\sqrt\eps)$,
$w_2=O(\eps^{\beta/2})=o(\eps)$.
Together with $G_{1,V}=2\eps F$, these prove the normal bounds
in \eqref{ent:threshold-normal} on every partial annulus.
In particular all energy, coefficient, force, chart and orientation
stops are strictly improved.

Along a rigid trace, division of
$\dot u_0=g_*+O(p)+O(p^{\beta-1})$ by
$\dot p=-a+O(p)$ gives
$u_{0,{\rm rig},\nu}=\gamma p/a+O(p^2)+O(p^\beta)$.
Detuning subtracts $k\sqrt\eps$ without changing that trace.
This proves the formula for $p_{\rm stop}$. At equal $p$, the
inverse weak-coordinate estimate gives on the annulus
\begin{equation}
 u-u_{\eta,{\rm rig},\nu}(p)
 =O(\eps^2+p^\beta\eps^{1/4})=o(\eps).
 \label{ent:physical-slip-match}
\end{equation}
Also $\dot u=g_*+O(\sqrt\eps)+O(\eps^{N_{\rm tail}})
<-\inf\gamma/2$. At the lower annular boundary the rigid slip
is at most $-k_-\sqrt\eps/2+O(\eps)$, contradicting the positive
slip stop and \eqref{ent:physical-slip-match}. The threshold is
therefore reached first and is crossed strictly once. Earlier slip
margins exclude a preceding crossing on this continuous passage.
Expansion at the same rigid zero gives more explicitly
\begin{equation}
 |R_e|\le C\{\eps^2+\eps^{\beta/2+1/4}
                  +\eps^{3/2}T+\eps^2T^2\}=o(\eps).
 \label{ent:threshold-remainder}
\end{equation}
If the profile tail is kept separately it adds
$C\eps^{N_{\rm tail}-1/4}=o(\eps)$. This proves the lemma.
\end{proof}

\subsection{The critical passage and its first force zero}

The threshold sets $\tau=(t-t_e)/\eps-T$ and $U=u/\eps$, so
$U(-T)=\gamma T$ exactly. Its normal error is
$E_{N,e}=O(\sqrt\eps)$, including $Z$, and
$(1+T)^2E_{N,e}\to0$. The power two matters because the saturated
critical propagator has a quadratic neutral part.

The limiting normal--slip equations are
\begin{equation}
 \begin{gathered}
 Y'=V,\qquad Z'=-Z-L_+,\qquad L=Z-Y,\\
 V'=C_G(1-\Phi_\theta(U))L_+,\\
 U'=d_\lambda+[-h+\mu A_G(1-\Phi_\theta(U))]L_+.
 \end{gathered}
 \label{ent:critical-field}
\end{equation}
with normalization
$(Y,V,Z,U-g_*\tau)\to(-2\ell,0,-\ell,0)$ as $\tau\to-\infty$.

\begin{lemma}[The normalized critical passage]
\label{lem:critical-pullback}
This normalization determines a unique solution $\overline W$ on
$\mathbb R$, with error $O(e^{\omega_c\tau})$ as $\tau\to-\infty$,
for common $\omega_c>0$. Its finite first force zero $\tau_r$ satisfies
\begin{equation}
 Y_r^0=Z_r^0<0,\qquad q_{\rm in}=V_r^0>0,\qquad U_r^0>0,
 \qquad q_{\rm in}+Y_r^0>0.
 \label{ent:critical-release-data}
\end{equation}
The solution, first-zero time and release data depend continuously
on $(\lambda,\theta)$; the four strict quantities have common
positive lower bounds on the fixed compact classes.
\end{lemma}
\begin{proof}
On a far left interval use positive contact. Define
$If(\tau)=\int_{-\infty}^{\tau}f(s)\,ds$ and
$Kf(\tau)=\int_{-\infty}^{\tau}e^{-2(\tau-s)}f(s)\,ds$.
The equations with their incoming limits are equivalent to
\[
 \begin{gathered}
 Y+2\ell=IV,\quad Z+\ell=KIV,\quad L=\ell+(K-\operatorname{Id})IV,
 \\
 U=g_*\tau-hI(L-\ell)+(\mu A_G/C_G)V,\\
 V=I[C_G(1-\Phi_\theta(U))L].
 \end{gathered}
\]
In $\|f\|_\omega=\sup_{\tau\le-T_0}e^{-\omega\tau}|f(\tau)|$,
where $0<\omega<\min(2,\alpha\inf\gamma/2)$, the operators
$I,K$ have norms $1/\omega,1/(2+\omega)$.
The first four formulas reconstruct the state boundedly from $V$.
On a fixed ball, sufficiently large common $T_0$ makes $L>0$ and
$U$ lie in the positive tail. Both $1-\Phi_\theta(U)$ and its
$U$ derivative are then bounded by $Ce^{\alpha\inf\gamma\tau}$.
The last formula is a self-map contraction, with constant
$Ce^{-\alpha\inf\gamma T_0}<1$. It gives the exponential tail.
Any solution with the stated incoming limits first has exponentially
small $V'$; integration from its zero incoming velocity gives these
same formulas and places it in the contraction ball. This proves
uniqueness. The locally Lipschitz positive-part field has at most
linear growth, so continuation is global.

The filter formula
$Z(\tau)=e^{-(\tau-s)}Z(s)-\int_s^\tau e^{-(\tau-\xi)}L_+(\xi)\,d\xi$
keeps $Z<0$. At $U=0$ its derivative is $d_\lambda+B_GL_+>0$, so $U$
stays positive. Until a first force zero, $V'>0$, hence $V>0$.
If force remained positive forever, $Y'=V$ would eventually make
$Y>0$, contrary to $Y<Z<0$. A finite first zero exists.
There $L'=-Z-L_+-V\le0$. If $L'=0$, differentiation from the
contact side gives $L''(\tau_r^-)=Z_r<0$, forcing $L<0$ just
before its first zero. Thus $L'=-(q_{\rm in}+Y_r^0)<0$.
Uniform contraction, continuous dependence on compact intervals,
and this transverse first zero prove continuous dependence of all
the data. Compactness gives the common strict margins.
\end{proof}

\begin{lemma}[Compactly supported variations of the critical passage]
\label{lem:profile-variation}
Fix $\lambda$ and an admissible $C^4$ baseline $\Phi^{\rm b}$.
Let $K\Subset(Q_\Phi,\infty)$ be a closed interval on which
$(\Phi^{\rm b})'$ has a positive lower bound, and put
\[
 E_K=\{\eta\in C^4(\R):\operatorname{supp}\eta\subset K\}.
\]
On a sufficiently small open ball about zero in $E_K$, equipped
with its $C^4$ norm, the normalized first-release time and state
for $\Phi^{\rm b}+\eta$ are Fr\'echet $C^1$.
For a direction $\psi\in E_K$, let $z$ be the fixed-time variation
of the contact extension, in state order $(Y,V,Z,U)$, and write
$q=q_{\rm in,\lambda}[\Phi^{\rm b}]=V(\tau_r)$ and
$y=Y(\tau_r)=Z(\tau_r)$. Then
\begin{equation}
 D\tau_r[\psi]=\frac{z_Z(\tau_r)-z_Y(\tau_r)}{q+y},
 \qquad Dq[\psi]=z_V(\tau_r).
 \label{ent:profile-event-derivative}
\end{equation}
Let $F_{\Phi^{\rm b}}$ be the contact extension of
\eqref{ent:critical-field}, obtained by replacing $L_+$ by $L$.
With $b=(0,C_G,0,\mu A_G)^T$,
$A=D_WF_{\Phi^{\rm b}}(\overline W)$ and
$\zeta'=-A^T\zeta$, $\zeta(\tau_r)=e_V$, one has
\begin{equation}
 Dq[\psi]=-\int_{-\infty}^{\tau_r}
           L(\tau)b^T\zeta(\tau)\psi(U(\tau))\,d\tau.
 \label{ent:profile-response}
\end{equation}
\end{lemma}
\begin{proof}
Choose a finite $\tau_a<\tau_r$ so negative that
$U(\tau)>\sup K+1$ and $L(\tau)>\ell/2$ for
$\tau\le\tau_a$. Every perturbation supported in $K$ vanishes
on this entire baseline tail. Attach the same value
$\overline W(\tau_a)$ to the perturbed forward initial-value
problem. The attached orbit has the prescribed normalization;
the tail uniqueness in Lemma~\ref{lem:critical-pullback} identifies
it with the perturbed pullback. Conversely, any normalized perturbed
orbit lies above $K$ sufficiently far in the past, where that same
uniqueness applies. Thus the incoming state at $\tau_a$ is exactly
independent of the shape.

Use the contact extension through the baseline first release.
Evaluation $(\eta,U)\mapsto\eta(U)$ is jointly $C^1$ on
$E_K\times\R$, so finite-time ODE dependence from the common
initial state is Fr\'echet $C^1$. At the event
$L'=-(q+y)<0$, and the implicit-function theorem gives a nearby
$C^1$ zero. The exact positive incoming tail, the positive minimum
of $L$ on the preceding compact interval, and strict decrease in
a small sign-changing event bracket exclude all earlier zeros.
The extension therefore agrees with the true positive-part solution
up to its first release.

The fixed-time variation satisfies
\[
 z'=Az-L\psi(U)b,\qquad z(\tau_a)=0.
\]
Differentiating $Z-Y=0$ gives the event-time formula. Since
$V'(\tau_r)=C_G[1-\Phi^{\rm b}(U(\tau_r))]L(\tau_r)=0$,
the moving-event contribution to $Dq$ vanishes.
For the fixed state event $H=Z-Y$ and output $J=V$, with no direct
time or profile dependence, the terminal covector is therefore
\[
 \nabla J-\frac{\nabla J\cdot F}{\nabla H\cdot F}\nabla H
 =e_V\ne0,
 \qquad \nabla H\cdot F=-(q+y)\ne0.
\]
Pairing with the terminal adjoint gives
$(\zeta^Tz)'=-L\psi(U)b^T\zeta$, hence
\eqref{ent:profile-response}. Its integrand vanishes before
$\tau_a$, so no integrability condition for the adjoint at negative
infinity is needed. Small perturbations in $E_K$ remain admissible
by the derivative margin on $K$, unchanged core and zero, and the
finite compact-support contribution to (F4). The weighted and
unweighted $C^4$ norms on $E_K$ are equivalent.
\end{proof}

\begin{theorem}[First release of the captured passage]
\label{thm:detuned-release}
Let $\sigma=\sqrt\eps$ and
\begin{equation}
 \mathcal D_\eps=(1+T)^2(E_{N,e}+e^{-\omega_c T})
                      +\sigma(1+T)^3+\eps(1+T)^4\longrightarrow0.
 \label{ent:critical-error}
\end{equation}
The contact passage from $p_i$ has a first release $t_r>t_e$.
Slip remains positive, and its complete release state satisfies
\begin{gather}
 t_r=t_e+\eps\{T+\tau_r+O(\mathcal D_\eps)\},\quad
 y_r=z_r=\eps^2\{Y_r^0+O(\mathcal D_\eps)\},
 \label{ent:release-time-memory}\\
 v_r=\eps\{q_{\rm in}+O(\mathcal D_\eps)\},\quad
 u_r=\eps\{U_r^0+O(\mathcal D_\eps)\},\quad
 \dot g_r=\eps\{q_{\rm in}+Y_r^0+O(\mathcal D_\eps)\}>c_0\eps,
 \label{ent:actual-release}\\
 |X_r-G|=O(\sigma),\qquad
 p_r/\sigma\to H:=ak/\gamma,\qquad b_r/\sigma\to-\ell H.
 \label{ent:release-slow}
\end{gather}
The leading normal--slip data depend only on $(\lambda,\theta)$.
\end{theorem}
\begin{proof}
The exact scaled equations have $Y',Z'$ as in
\eqref{ent:critical-field} and
\[
 V'=b+[p+\mu B(1-\Phi_\theta(U))]L_+,\qquad
 U'=a_T+[B-\mu A+\mu A(1-\Phi_\theta(U))]L_+.
\]
Stop while $(Y,V,Z)$ stays in a fixed enlarged bounded neighborhood
of the reference normal range, $|U-\overline U(\tau)|\le1$,
and the mechanical state stays in a common chart. On $[-T,-T_0]$,
also stop at the boundary of the convex positive-contact tail
neighborhood used below. On the tail,
$\overline U(\tau)=g_*\tau+O(e^{\rho\tau})$, so this tube and
every joining segment satisfy $U\ge-\gamma\tau-C$.
On every partial interval of length $T+C$, the mechanical field is bounded and
$|X-G|\le CM_\eps$, $M_\eps=\sigma+\eps(1+T)$.
Subtracting \eqref{ent:critical-field} gives zero residual in $Y,Z$
and residuals
\[
 b+pL_++(\mu B-C_G)(1-\Phi_\theta(U))L_+,\qquad
 (a_T-d_\lambda)+(B-\mu A+h)L_++\mu(A-A_G)(1-\Phi_\theta(U))L_+,
\]
of size $CM_\eps$. The term $b+pL_+$ is included.

On $[-T,-T_0]$, the saturated linear part is
\[
 M_\infty=\begin{pmatrix}0&1&0&0\\0&0&0&0\\1&0&-2&0\\h&0&-h&0\end{pmatrix}.
\]
Its homogeneous solutions have constant $V$, affine $Y$, a decaying
$Z-Y/2+V/4$, and at most quadratic $U$; hence
$\|e^{M_\infty s}\|\le C(1+s)^2$ for $s\ge0$.
In a convex positive-contact tail neighborhood the remaining
nonlinearity is
$(1-\Phi_\theta(U))L(0,C_G,0,\mu A_G)^T$ and its derivative is
$Ce^{\omega_c\tau}$. Decrease $\omega_c$ if needed so the same
exponent applies here and to the pullback tail. This follows from
(F4) on every joining segment.
Thus the profile tail has all needed moments through degree two;
no slow-time derivative of its composition is used.
For $D=|W_\eps-\overline W|$, Duhamel gives on each partial interval
\[
 \begin{aligned}
 D(\tau)&\le C(1+\tau+T)^2a_\eps
 +C\int_{-T}^{\tau}(1+\tau-s)^2e^{\omega_cs}D(s)\,ds,\\
 a_\eps&\le C\{E_{N,e}+e^{-\omega_c T}
                      +\sigma(1+T)+\eps(1+T)^2\}.
 \end{aligned}
\]
Divide by $(1+\tau+T)^2$ and take the supremum. The integral
coefficient is at most
$Ce^{-\omega_c T_0}\int_0^\infty(1+s)^2e^{-\omega_cs}\,ds<1/2$
after increasing $T_0$. Therefore
$\sup_{[-T,-T_0]}D\le C\mathcal D_\eps$.
This improves the scaled-state, positive-force and chart stops.
On a common fixed interval from $-T_0$ past all reference releases,
the positive-part field is uniformly Lipschitz. Stopped Gronwall
gives the same bound and improves the remaining state stop.

The reference signed quantity $L=Z-Y$ is positive before a common small bracket of
$\tau_r$, has opposite endpoint signs there, and has a strictly
negative derivative in the bracket. The exact identity
$L'=-Z-L_+-V$ transfers this derivative margin and the signs to
the actual trajectory. It therefore has exactly one zero there,
at $\tau_r+O(\mathcal D_\eps)$, and no preceding zero.
The threshold lemma supplies contact from $p_i$ to $t_e$.
At physical $u=0$, $\dot u=a_T+B\lambda_N>0$ on the chart,
so the initial positive slip persists. At the zero,
$\dot g=\eps(V+Z)$; the data and strict opening follow.
Finally the physical duration is $O(\eps(1+T))=o(\sigma)$.
Integration of $\dot p=-a+O(M_\eps)$ gives
\[
 p_r=p_e-a\eps(T+\tau_r)
       +O(\eps\mathcal D_\eps+\eps(1+T)M_\eps).
\]
The $a\eps T$ term cancels before using
\eqref{ent:threshold-remainder}; the $b$ change is
$O(\eps(1+T))$. This proves \eqref{ent:release-slow}.
Uniqueness of the normalized critical solution proves the asserted
dependence of the leading data.
\end{proof}
